% Part: counterfactuals
% Chapter: introduction
% Section: strict-conditional

\documentclass[../../../include/open-logic-section]{subfiles}

\begin{document}

\olfileid{cnt}{int}{str}

\olsection{कठोर अभिव्यंजन}

लुईस यांनी \emph{कठोर अभिव्यंजन}~$\strictif$
हा कारक मांडला आणि वास्तविक अभिव्यंजनाऐवजी
तो तार्किक निष्पन्नतेशी जुळतो, असा युक्तिवाद
केला. अवश्यतावाचक मोडल तर्कशास्त्रात
$!A \strictif !B$ ची व्याख्या
$\Box(!A \lif !B)$ अशी करता येते.
म्हणून एखाद्या जगात कठोर अभिव्यंजन
सत्य असते तेव्हा आणि केवळ तेव्हाच संबंधित
वास्तविक अभिव्यंजन अवश्य सत्य असते.

वास्तविक अभिव्यंजनाच्या विरोधाभासांच्या
बाबतीत कठोर अभिव्यंजन कसे वागते?
पूर्वांग असत्य असलेले कठोर अभिव्यंजन,
तसेच उत्तरांग सत्य असलेले कठोर अभिव्यंजन,
सत्यही असू शकते किंवा असत्यही.
शिवाय $(!A \strictif !B) \lor (!B \strictif !A)$
वैध नाही. $!A \strictif !B$
हे कठोर अभिव्यंजन
$\lnot !A \lor !B$ शीही
सममूल्य नाही; त्यामुळे ते
सत्यता-फलनात्मक नाही.

S5 च्या संदर्भात पुढील संबंध मिळतात:
\begin{align}
  !A \strictif !B & \Entails \lnot !A \lor !B
  \text{ पण:}\\
  \lnot !A \lor !B & \Entails/ !A \strictif !B\\
  !B & \Entails/ !A \strictif !B\\
  \lnot !A & \Entails/ !A \strictif !B\\
  \lnot(!A \strictif !B) & \Entails/ !A \land \lnot !B
  \text{ पण:}\\
  !A \land \lnot !B & \Entails \lnot(!A \strictif !B)
  \intertext{तरीही कठोर अभिव्यंजनामुळे
    विध्यात्मक अनुमान लागू होते:}
  !A, !A \strictif !B & \Entails !B
  \intertext{कठोर अभिव्यंजन
    निष्कर्षांचे संयुग्मीकरण करते:}
  !A \strictif !B,  !A \strictif !C & \Entails !A \strictif (!B \land !C)
  \intertext{कठोर अभिव्यंजनासाठी
    पूर्वांग अधिक बलवान करता येते:}
  !A \strictif !B & \Entails (!A \land !C) \strictif !B
  \intertext{कठोर अभिव्यंजन
    संक्रमकही आहे:}
  !A \strictif !B, !B \strictif !C & \Entails !A \strictif !C
  \intertext{शेवटी, कठोर अभिव्यंजन
    त्याच्या प्रतिपरिवर्तित विधानाशी
    सममूल्य आहे:}
  !A \strictif !B & \Entails \lnot !B \strictif \lnot !A\\
  \lnot !B \strictif \lnot !A & \Entails !A \strictif !B
\end{align}

\begin{prob}
  कठोर अभिव्यंजनासाठी लागू न होणाऱ्या
  तार्किक निष्पन्नता-संबंधांना
  \Log{S5} मधील प्रतिदृष्टांत द्या;
  म्हणजे पुढील प्रत्येकासाठी:
  \begin{enumerate}
  \item $\lnot p \Entails/ \Box(p \lif q)$
  \item $q \Entails/ \Box(p \lif q)$
  \item $\lnot\Box(p \lif q) \Entails/ p \land \lnot q$
  \item $\Entails/ \Box(p \lif q) \lor \Box(q \lif p)$
  \end{enumerate}
\end{prob}

\begin{prob}
  वैध तार्किक निष्पन्नता-संबंध
  कठोर अभिव्यंजनासाठी लागू होतात
  हे दाखवण्यासाठी पुढील
  विधानांचे \Log{S5} मध्ये
  पुरावे द्या:
  \begin{enumerate}
  \item  $\Box(!A \lif !B) \Entails \lnot !A \lor !B$
  \item $!A \land \lnot !B \Entails \lnot\Box(!A \lif !B)$
  \item $!A, \Box(!A \lif !B) \Entails !B$
  \item $\Box(!A \lif !B), \Box(!A \lif !C) \Entails \Box(!A \lif (!B
    \land !C))$
  \item $\Box(!A \lif !B) \Entails \Box((!A \land !C) \lif !B)$
  \item $\Box(!A \lif !B), \Box(!B \lif !C) \Entails \Box(!A
    \lif !C)$
  \item $\Box(!A \lif !B) \Entails \Box(\lnot !B \lif \lnot !A)$
  \item $\Box(\lnot !B \lif \lnot !A) \Entails \Box(!A \lif !B)$
  \end{enumerate}
  \end{prob}

तरी कठोर अभिव्यंजनाचे स्वतःचेही
‘‘विरोधाभास’’ आहेत. पूर्वांग असत्य
असेल किंवा उत्तरांग सत्य असेल,
तर वास्तविक अभिव्यंजन सत्य
असते. त्याचप्रमाणे, पूर्वांग
\emph{अवश्य} असत्य असेल
किंवा उत्तरांग अवश्य सत्य असेल,
तर कठोर अभिव्यंजन सत्य असते.
S5 मध्ये कोणतेही सत्य कठोर
अभिव्यंजन अवश्य सत्य असते
आणि कोणतेही असत्य कठोर
अभिव्यंजन अवश्य असत्य असते.
दुसऱ्या शब्दांत,
\begin{align}
  \Box \lnot !A & \Entails !A \strictif !B\\
  \Box !B & \Entails !A \strictif !B\\
  !A \strictif !B& \Entails \Box(!A \strictif !B)\\
  \lnot(!A \strictif !B) & \Entails \Box\lnot(!A \strictif !B)
\end{align}
$\strictif$ चा अर्थ ‘‘तार्किकरीत्या
निष्पन्न होते’’ असा घेतल्यास
यात अडचण नाही. शेवटी,
तार्किक निष्पन्नतेचे संबंध ही
गणिती तथ्ये आहेत; म्हणून
ती नैमित्तिक असू शकत नाहीत.
पण ‘‘जर \dots तर \dots’’
याचा अर्थ पकडणारा तार्किक
कारक म्हणून $\strictif$
वापरायचा असेल, तर या
गोष्टी प्रश्न निर्माण करतात,
विशेषतः शेवटच्या दोन.
कारण काही ‘‘जर \dots तर \dots’’
विधाने नैमित्तिकरीत्या
सत्य किंवा असत्य असतात;
प्रत्यक्षात ती सहसा
ना अवश्य सत्य असतात,
ना अशक्य.

\begin{prob}
  पुढील विधानांचे \Log{S5}
  मध्ये पुरावे द्या:
  \begin{enumerate}
    \item  $\Box \lnot !A \Entails !A \strictif !B$
    \item $!A \strictif !B \Entails \Box(!A \strictif !B)$
    \item $\lnot(!A \strictif !B)  \Entails \Box\lnot(!A \strictif !B)$
  \end{enumerate}
  त्यासाठी $\strictif$ ची
  व्याख्या वापरा.
\end{prob}

\end{document}
